\chapter{Introduction and Reading Guide}
\label{ch:introduction}

This document develops one construction and follows it as far as it goes. A population of agents is
carried by an associated statistical bundle over a contextual base; a single fixed normalized law
generates their latent states and observations; a single correlated recognition law approximates the
resulting posterior; and one exact identity relates the two. Everything after that is either a
consequence of that identity, a declared restriction of the family it is evaluated on, or a
statement about what happens to the whole structure when the population is described at a coarser
resolution.

The last of these is the reason the document exists. Coarse-graining a population of inferring
agents is a renormalization, and asking which coarse descriptions the exact bound licenses turns out
to be a sharper question than it first appears. Some operations that are routinely called
coarse-graining move the evidence, some move only the bound, and some are not admissible at all.
Separating them is most of the work, and the separation has consequences the informal picture does
not anticipate.

\section*{The epistemic contract}

The single most important convention in this document is that every non-trivial claim carries a
visible status, and that the statuses are not interchangeable. They appear as a bracketed tag
immediately after the statement they govern, and they are also named in the surrounding prose so
that the discipline survives being read aloud.

\begin{sciencetable}{Status tags used throughout. A claim without a tag is an error, and a
conjecture written in the grammar of a theorem is a worse error.}
\label{tab:status-taxonomy}
\begin{tabularx}{\textwidth}{>{\raggedright\arraybackslash}p{0.19\textwidth}>{\raggedright\arraybackslash}X}
\toprule
\textbf{Tag} & \textbf{What it promises the reader} \\
\midrule
\status{ESTABLISHED} & Proved here, or a standard result cited to a source that has been checked. \\
\status{DEFINITION} & A declared type, construction, or convention. Nothing is being proved and the text says so. \\
\status{HYPOTHESIS} & A restriction the development adopts by choice. What it excludes, and where it is used, are stated. \\
\status{CONJECTURE} & Believed and precisely stated, but not proved. Stated sharply enough to be attacked. \\
\status{NUMERICAL} & Supported by computation only, with its measurement, seed, and control reported. Computation is not proof. \\
\status{OPEN} & Unsettled. What would settle it, and what obstructs it, are named. \\
\status{NOT-CLAIMED} & A statement the development deliberately declines to make. Declining is not refuting, and the text says which it is doing. \\
\bottomrule
\end{tabularx}
\end{sciencetable}

Two habits follow from the table and are worth stating explicitly, because they are the ones most
easily lost.

A hand-wave is treated as a debt rather than a style. Any step made by gesture rather than argument
is either discharged, by supplying the proof or computing the constant or naming the exact condition,
or else converted into a declared gap that states what holds, under what hypotheses, what remains
unproved, and what would close it. Weakening the verb is not a permitted third option: a sentence
reporting that something presumably converges is worse than one reporting that convergence is open,
because the first conceals the debt inside a construction that reads like a result.

Declining a claim and refuting one are kept apart. Both occur here. Chapter~\ref{ch:obstructions}
records constructions that fail, and several chapters record adjacent statements the development
deliberately does not make. Conflating the two would misrepresent the theory in opposite directions,
and each chapter closes with a register listing its own results and their statuses so that the
distinction can be audited rather than trusted.

\section*{What is established, and what is not}

It is worth saying at the outset, in compressed form, where this development actually stands. The
foundations are secure: the typed construction, the fixed normalized generative law with its gauge
covariance, and the exact evidence identity are proved under hypotheses that are stated and used.
The Gaussian realization is likewise secure, with one qualification that matters and is easy to miss,
namely that the interaction family on which the entire second half operates is a \emph{declared
subfamily} and not a consequence of the linear-Gaussian model. Chapter~\ref{ch:gaussian} exhibits
the parameter choices under which the assumed form fails.

The coarse-graining results are proved, and the one most likely to be misread is negative. Closure of
the interaction family under aggregation does hold, exactly, but it does not characterize the family:
a one-parameter continuum of alternatives is equally closed, and what actually selects the form is a
symmetry rather than a renormalization property. The coarse operator itself is not new, and is cited
to the numerical-analysis literature rather than claimed.

The renormalization results are the least settled, and the document is arranged so that this is
impossible to overlook. The rescaling that turns aggregation into a flow is declared, not derived.
The scale-free fixed point is proved invariant, which is strictly weaker than the attraction that
universality would require, and attraction is open. Evidence that the flow moves toward that fixed
point is numerical, obtained against matched controls, and is reported as numerical. The
identification of a fixed point with a physical law, which is what makes the construction interesting
to anyone who cares about where physical theory comes from, is a declaration and is marked as one
everywhere it appears.

\section*{Scope}

The construction here is general. It is not a commentary on, a defense of, or a crosswalk to any
particular earlier formulation, and it borrows no term structure, no engineered objective, and no
implementation detail from one. Where a prior construction is refuted, in
Chapter~\ref{ch:obstructions}, the refutation is stated as a general structural result about the
class of constructions it belongs to, because that is the form in which it is actually true and the
form in which it is useful to a reader who has never seen the specific case.

Two chapters are load-bearing for readers in a hurry. Chapter~\ref{ch:elbo} contains the identity
everything rests on. Chapter~\ref{ch:coarsegraining} contains the operations, their costs, and the
result that closure is not a characterization. A reader interested only in the renormalization may
begin at Chapter~\ref{ch:gaussian} and refer back, at the cost of taking the typing on trust.
